% 03_system_model.tex  --  WRITTEN
\section{System Model and Problem Formulation}
\label{sec:system_model}

We consider a single SU opportunistically accessing a
wideband spectrum of bandwidth $\Wband$ centered at $\fc$, which is
shared by an \emph{unknown and time-varying} number of PU
emissions. Such a setting arises in shared bands such as CBRS, where
incumbents select center frequencies and bandwidths adaptively rather than on
a fixed channel grid~\cite{sohul2015cbrs,xing2007dsa}. Unlike the
fixed-channel models prevalent in the DRL-based DSA literature, we make
\emph{no} assumption that PU activity is aligned to a predefined channel
grid: each PU emission occupies an arbitrary center frequency, bandwidth,
start time, and duration, exactly the regime targeted by TF
localization~\cite{prasad2020,zhao2025trtfl_ssl}. The notation introduced
below is summarized in Table~\ref{tab:notation}.

\begin{table}[t]
\centering
\caption{Summary of Notation}
\label{tab:notation}
\begin{tabular}{ll}
\toprule
Symbol & Meaning \\
\midrule
$\spec\in\R_{\ge0}^{\Nf\times\Nt}$ & Magnitude spectrogram (frame) \\
$\cell$ & TF cell (freq.\ bin $m$, time bin $n$) \\
$\boxk_k=(\bfc_k,\bBW_k,\bts_k,\bdur_k)$ & $k$-th PU emission box \\
$\occmap\in\{0,1\}^{\Nf\times\Nt}$ & True occupancy map \\
$\belmap\in[0,1]^{\Nf\times\Nt}$ & Estimated occupancy belief map \\
$\Pmd,\Pfa$ & Cell-level miss/false-alarm prob.\ \\
$\thr$ & Belief decision threshold \\
$\act=(\afc,\aBW,\ats,\adur,\pwr)$ & SU access action (resource block) \\
$\dmin,\pmax$ & Min.\ bandwidth demand, power budget \\
$\state_t,\obs_t$ & Access state, observation \\
$\rew,\disc,\pol$ & Reward, discount, policy \\
$\Pcol,\Pcolmax$ & PU collision prob.\ and its cap \\
$\thru,\thrumax$ & Achievable / oracle throughput \\
\bottomrule
\end{tabular}
\end{table}

\subsection{Wideband Signal and Spectrogram Model}
\label{subsec:signal}
The complex baseband signal observed by the wideband sensor over one frame
is
\begin{equation}
\rx(t) \;=\; \sum_{k=1}^{\Npu} c_k\!\left(\sutx_k(t)\right) \;+\; \noise(t),
\qquad 0 \le t < \bigT_{\mathrm{f}},
\label{eq:rx}
\end{equation}
where $\sutx_k(t)$ is the $k$-th PU emission, $c_k(\cdot)$ denotes the
corresponding channel (carrying carrier offset, bandwidth-limiting, and
fading), $\noise(t)$ is additive white Gaussian noise, $\Npu$ is the
(random, unknown) number of active emissions in the frame, and
$\bigT_{\mathrm{f}}$ is the frame length. The signal is sampled at rate
$\fs \ge \Wband$.

Applying the short-time Fourier transform (STFT) operator $\STFT$ yields a
single-resolution TF representation; we work with the
magnitude spectrogram
\begin{equation}
\spec \;=\; \bigl|\STFT\{\rx\}\bigr| \;\in\; \R_{\ge 0}^{\,\Nf \times \Nt},
\label{eq:stft}
\end{equation}
discretized into $\Nf$ frequency bins and $\Nt$ time bins. We index an
individual TF \emph{cell} by $\cell$, $m\in\{1,\dots,\Nf\}$,
$n\in\{1,\dots,\Nt\}$. The STFT is preferred over variable-resolution
transforms (the continuous wavelet transform or constant-$Q$ transform)
because spectrum monitoring weights all frequencies equally~\cite{prasad2020}.

\subsection{Continuous Occupancy Representation}
\label{subsec:occupancy}
Each PU emission $k$ is described by a continuous TF \emph{box}
\begin{equation}
\boxk_k \;=\; \bigl(\bfc_k,\, \bBW_k,\, \bts_k,\, \bdur_k\bigr),
\label{eq:box}
\end{equation}
i.e.,\ center frequency, bandwidth, start time, and duration. The box maps to
a region $\mathcal{B}_k \subset \{1,\dots,\Nf\}\times\{1,\dots,\Nt\}$ of the
discretized TF plane. The \emph{true occupancy map} is the binary matrix
$\occmap\in\{0,1\}^{\Nf\times\Nt}$ with
\begin{equation}
\occ_{m,n} \;=\;
\Ind\!\Bigl\{\, \cell \in \textstyle\bigcup_{k=1}^{\Npu}\mathcal{B}_k \,\Bigr\},
\label{eq:occupancy}
\end{equation}
where $\occ_{m,n}=1$ marks a cell occupied by at least one PU. Crucially,
$\occmap$ is \emph{not} block-structured along a fixed channelization;
recovering it is exactly the TF-localization problem.

We clarify the sense in which the model is ``grid-free.'' The PU
\emph{emission parameters} \eqref{eq:box} are continuous; they are mapped to
an occupancy map discretized at the \emph{sensing resolution} $\Nf\times\Nt$
of the STFT. The point is not the absence of any discretization---a finite
spectrogram is necessarily discretized---but the absence of a fixed
\emph{communication channel grid} to which PU activity is assumed aligned.
Emissions may occupy arbitrary cells at sensing resolution, straddling any
notional channel boundary, which is precisely what defeats fixed-channel DSA
models and motivates the localization-driven treatment.

\subsection{Sensing Model: From Spectrogram to Belief Map}
\label{subsec:sensing}
The sensing module is a learned mapping $\phi_{\mathrm{sen}}$ that estimates,
for each cell, the posterior probability of PU occupancy:
\begin{equation}
\belmap = \phi_{\mathrm{sen}}(\spec) \in [0,1]^{\Nf\times\Nt},
\quad
\bel_{m,n} \approx \Prob(\occ_{m,n}{=}1 \mid \spec).
\label{eq:belief}
\end{equation}
Concretely, $\bel_{m,n}$ is the \emph{detection confidence} of the
time--frequency localizer rasterized onto the cell grid: each estimated
emission box carries a confidence score, and $\bel_{m,n}$ is the (edge-feathered)
score of the box covering cell $(m,n)$, or a low background value where no box
is predicted. We deliberately retain this \emph{soft} confidence $\belmap$
rather than applying NMS to obtain hard boxes; this
preserves the sensing uncertainty that the access stage will exploit
(cf.\ Section~\ref{sec:framework}). The reliability of $\phi_{\mathrm{sen}}$
is characterized at the cell level by the miss-detection and false-alarm
probabilities
\begin{align}
\Pmd &\defeq \Prob\!\left(\bel_{m,n} < \thr \;\middle|\; \occ_{m,n}=1\right),
\label{eq:pmd}\\
\Pfa &\defeq \Prob\!\left(\bel_{m,n} \ge \thr \;\middle|\; \occ_{m,n}=0\right),
\label{eq:pfa}
\end{align}
where $\thr\in[0,1]$ is a decision threshold applied to the localizer's
confidence $\bel_{m,n}$. Two distinct factors shape $(\Pmd,\Pfa)$: the
\emph{confidence threshold} $\thr$ (where we cut the soft score) and the
\emph{geometric accuracy} of the box (how well its frequency support matches
the true emission, quantified later by the intersection-over-union (IoU)
error $\eloc$). The
operating-point analysis of Theorem~\ref{thm:tradeoff} optimizes over $\thr$,
while the localization-error model of Assumption~\ref{as:localization}
captures the geometric contribution; their separation is made explicit in the
bound \eqref{eq:loc_to_cell}. Both $\Pmd$ and $\Pfa$ are functions of $\thr$;
this dependence is the lever in the trade-off of
Theorem~\ref{thm:tradeoff} (Section~\ref{sec:theory}).

\subsection{Access Model: Continuous Resource Selection}
\label{subsec:access}
The SU operates a sense-then-access frame: the sensing phase estimates the
occupied frequency ranges, and the access phase of the \emph{same} frame
transmits in the ranges it deems idle. Because both phases share the frame,
temporal alignment is given, and the access decision turns on \emph{which
frequencies} are free---making frequency-localization accuracy, rather than
temporal extent, the quantity that governs access performance (formalized in
Section~\ref{sec:theory}). At each frame the SU chooses a continuous TF
resource block together with a transmit power level,
\begin{equation}
\act \;=\; \bigl(\afc,\, \aBW,\, \ats,\, \adur,\, \pwr\bigr),
\label{eq:action}
\end{equation}
subject to the bandwidth demand, power, and PU-protection constraints
\begin{equation}
\aBW \ge \dmin, \quad 0 \le \pwr \le \pmax, \quad
\Prob(\text{collision}) \le \Pcolmax .
\label{eq:constraints}
\end{equation}
A \emph{collision} occurs when the SU block overlaps an occupied cell:
\begin{equation}
\mathrm{Col}(\act) \defeq
\Ind\!\Bigl\{\, \exists\,\cell \in \mathcal{A}(\act):\; \occ_{m,n}=1 \,\Bigr\},
\label{eq:collision}
\end{equation}
where $\mathcal{A}(\act)$ is the set of TF cells covered by action $\act$.
Transmission succeeds (yielding an acknowledgment, ACK) iff $\mathrm{Col}(\act)=0$ and the
allocated bandwidth meets the demand.

\subsection{Partially Observed Control Formulation}
\label{subsec:pomdp}
The SU never observes the true occupancy $\occmap$ directly; it observes only
the wideband spectrogram $\spec_t$ and, after transmitting, a one-bit
ACK or negative-acknowledgment (NACK). The control problem is therefore
partially observed. We make the
roles explicit to avoid conflating the detector output with a Bayesian
belief.

\paragraph*{Hidden state} The hidden state is the true PU occupancy together
with the underlying emission parameters,
$\hidden_t=\bigl(\occmap_t,\{\boxk\}\bigr)$, evolving as a
(piecewise-stationary) Markov process.

\paragraph*{Observation} The raw observation is the spectrogram $\spec_t$ and
the ACK/NACK $\mathrm{ACK}_t\in\{0,1\}$. The sensing tier maps $\spec_t$ to
the per-cell soft occupancy estimate $\belmap_t\in[0,1]^{\Nf\times\Nt}$ of
\eqref{eq:belief}. We stress that $\belmap_t$ is a learned \emph{soft
occupancy observation feature}, not the exact Bayesian posterior
$\Prob(\hidden_t\mid\text{history})$; it plays the role of an
\emph{approximate sufficient statistic} of the observation history for the
access decision.

\paragraph*{Access state} The access policy acts on the recent history of
this feature,
\begin{equation}
\state_t \;=\; \bigl(\belmap_{t-\histlen+1},\dots,\belmap_{t}\bigr)
\;\in\; [0,1]^{\histlen \times \Nf \times \Nt},
\label{eq:state}
\end{equation}
stacked over a short window to capture temporal correlation of PU activity.
This continuous, uncertainty-bearing state is the central departure from
prior DRL-DSA work, whose state is a binary fixed-channel occupancy vector.
Because $\belmap$ approximates rather than equals the exact posterior, the
formulation is a \emph{partially observed control problem with a learned
observation feature} rather than a belief-MDP in the strict Bayesian sense;
the theoretical claims of Section~\ref{sec:theory} are stated against this
representation and do not presume exact belief recursion.

\paragraph*{Reward} We use an uncertainty-aware reward that rewards
throughput, penalizes PU collisions, and discourages transmitting into
low-confidence regions:
\begin{equation}
\rew(\state_t,\act_t) \;=\;
\underbrace{\thru(\act_t)}_{\text{throughput}}
\;-\; \lambda_{\mathrm{c}}\,\mathrm{Col}(\act_t)
\;-\; \lambda_{\mathrm{u}}\,\mathcal{U}\!\left(\belmap_t,\act_t\right),
\label{eq:reward}
\end{equation}
where $\mathcal{U}(\cdot)$ measures the residual occupancy uncertainty over
the chosen block and $\lambda_{\mathrm{c}},\lambda_{\mathrm{u}}\ge 0$ are
weights. The connection of the $\lambda_{\mathrm{u}}$ term to a
risk-sensitive objective is discussed in Section~\ref{sec:framework}.

\paragraph*{Objective} The SU seeks a policy
$\pol:\mathcal{S}\to\mathcal{A}$ maximizing the expected discounted return
subject to the PU-protection constraint,
\begin{equation}
\begin{aligned}
\polopt = \argmax_{\pol}\;
&\E_{\pol}\!\left[\sum_{t=0}^{\infty}\disc^{t}
\rew(\state_t,\act_t)\right]\\[-1pt]
\text{s.t.}\quad
&\E_{\pol}\!\left[\mathrm{Col}(\act_t)\right] \le \Pcolmax .
\end{aligned}
\label{eq:objective}
\end{equation}

\begin{remark}
Setting $\belmap$ to a binary fixed-channel vector and restricting $\act$ to
a single grid channel recovers the standard formulation of
\cite{wang2018,zhong2019} as a special case. Proposition~\ref{thm:sufficiency}
formalizes that this restriction can be strictly suboptimal under grid misalignment.
\end{remark}
